\newif\ifglobal

%%%%%%%%%%%%%%%%%%%%%%%
%%% Compile to view %%%
%%%%%%%%%%%%%%%%%%%%%%%

%%%%%%%%%%%%%%%%%%%%%%%%%%%%%%%%%%%%%%%%%%%%%%%%%%%
%%% Readme:                                     %%%
%%% https://www.overleaf.com/read/bwdjvfjksvjy/ %%%
%%%%%%%%%%%%%%%%%%%%%%%%%%%%%%%%%%%%%%%%%%%%%%%%%%%

% >>> M?: marks a module setting number or important notice

% >>> M4: Set {./relative-path} to external files for this module <<<
\def \modulefiles {./06-RESCLK}
% To add an external file, use:
% (Replace '---' with actual filename)
% '\input{\modulefiles/tables/---.tex}' for tables
% '\includegraphics{\modulefiles/figures/---}' for figures
% '\subfile{\modulefiles/---}' for register files

% >>> M1: This module is in global project? <<<
% 'YES' - uncomment; 'NO' - comment out
\globaltrue

\ifglobal
    \documentclass[main\_\_CN.tex]{subfiles}

\else
    % Font "FandolSong-Regular" used by ctex does not have GSUB table.
    % It causes warnings that do not affect compilation and can be ignored.
    \PassOptionsToPackage{quiet}{fontspec}

    \documentclass[openany, 10pt]{book}

    % >>> M2: In ./00-shared/preamble-trm-module, also find
    %         settings for visibility of todo notes, watermark <<<
    \usepackage{preamble-trm-module}


    % Enable Chinese typesetting
    \usepackage{ctex}

    % >>> M3: Temporary labels in Chinese
    %         you can add your own labels there <<<
    \myexternaldocument{./00-shared/temp-labels-cn}

\fi %\ifglobal


\setcounter{secnumdepth}{4}


% >>> M5: If this module needs any updates in
% './00-shared/preamble-trm-module.sty'
% (include additional package, etc.), contact your module PM
% or create the file './00-shared/changelog.tex' make the
% necessary updates yourself and document them in this file <<<

% Make text left-aligned and allow latex to break pages naturally, comment out for full justification
\raggedright
\raggedbottom


\begin{document}


% Sets language into which pre-defined bits of text are translated
\selectlanguage{chinese}
\setchineseformatting

% Do not delete this block
\ifglobal
% Add the following front matter if
% this module is outside of global project
\else
    \InputIfFileExists{readme}{}{}
    \listoftodos
    {\let\clearpage\relax \listofdonetodos}
    \tableofcontents
    %\listoftables
    %\listoffigures
\fi



%%%%%%%%%%%%%%%%%%%%%%%%%
%%% TRM Module Begins %%%
%%%%%%%%%%%%%%%%%%%%%%%%%

\hypertarget{resclk}{}
\chapter{复位和时钟}
\label{mod:resclk}


\section{复位}
\subsection{概述}
系统提供四种级别的复位方式，分别是 CPU 复位、内核复位、系统复位和芯片复位。

除芯片复位外其他复位方式不影响片上内存存储的数据。
图 \ref{figure:Reset} 展示了整个芯片系统的结构以及四种复位等级。

\begin{figure}[H]
    \centering
    \includegraphics[width=0.8\textwidth]{\modulefiles/figures/Reset.eps}
    \caption{四种复位等级}
    \label{figure:Reset}
\end{figure}

\begin{itemize}
    \item CPU 复位：只复位 CPU core，复位释放后，程序将从 CPU reset vector 开始执行；
    \item 内核复位：复位除 RTC 以外的其他数字系统，包括 CPU、外设、Wi-Fi 及数字 GPIO；
    \item 系统复位：复位包括 RTC 在内的整个数字系统；
    \item 芯片复位：复位整个芯片。

\end{itemize}

\subsection{复位源}
上述任一复位源产生时，CPU 均将立刻复位。复位释放后，CPU 可以通过读取寄存器 \hyperref[fielddesc:RTCCNTLRESETCAUSEPROCPU]{RTC\_CNTL\_RESET\_CAUSE\_PROCPU} 来获取复位源。

表 \ref{table:resetreasons} 列出了从这一寄存器中可能读出的复位源。

\begin{longtable}[c]{ |p{0.8cm}|p{3.5cm}|p{3.5cm}|p{8.2cm}| }
\caption{复位源}
    \label{table:resetreasons}  \\
        \hline
        \rowcolor{lightgray}
        编码 & 复位源 & 复位方式 & 注释\\
        \hline
        \endfirsthead
        \hline
       \rowcolor{lightgray}
        编码 & 复位源 & 复位方式 & 注释 \\
        \hline
        \endhead
        \endfoot
        \hexprefix{}01 & 芯片复位 & 芯片复位 &  见表下方芯片复位的触发源\\
        \hline
        \hexprefix{}0F & 欠压系统复位 & 系统复位 & 欠压检测器触发的系统复位\\
        \hline
        \hexprefix{}10 & RWDT 系统复位 & 系统复位 & 详见章节 \ref{mod:wdt} \textit{\nameref{mod:wdt}}\\
        \hline
        \hexprefix{}12 & 超级看门狗复位 & 系统复位 & 详见章节 \ref{mod:wdt} \textit{\nameref{mod:wdt}} \\
        \hline
        \hexprefix{}13 & GLITCH 复位 & 系统复位 & - \\
        \hline
        \hexprefix{}03 & 软件系统复位 & 内核复位 & 配置 \hyperref[fielddesc:RTCCNTLSWSYSRST]{RTC\_CNTL\_SW\_SYS\_RST} 寄存器触发\\
        \hline
        \hexprefix{}05 & Deep-sleep 复位 & 内核复位 & 详见章节 \ref{mod:lowpowm} \textit{\nameref{mod:lowpowm}} \\
        \hline
        \hexprefix{}07 & MWDT0 全局复位 & 内核复位 & 详见章节 \ref{mod:wdt} \textit{\nameref{mod:wdt}}\\
        \hline
        \hexprefix{}08 & MWDT1 全局复位 & 内核复位 & 详见章节 \ref{mod:wdt} \textit{\nameref{mod:wdt}}\\
        \hline
        \hexprefix{}09 & RWDT 内核复位 & 内核复位 & 详见章节 \ref{mod:wdt} \textit{\nameref{mod:wdt}}\\
        \hline
        \hexprefix{}14 & eFuse 复位 & 内核复位 & eFuse CRC 校验错误触发复位 \\
        \hline
        \hexprefix{}0B & MWDT0 CPU 复位 & CPU 复位 & 详见章节\ref{mod:wdt} \textit{\nameref{mod:wdt}}\\
        \hline
         \hexprefix{}0C & 软件 CPU 复位 & CPU 复位 & 配置 \hyperref[fielddesc:RTCCNTLSWPROCPURST]{RTC\_CNTL\_SW\_PROCPU\_RST} 寄存器触发\\
        \hline
        \hexprefix{}0D & RWDT CPU 复位 & CPU 复位 & 详见章节 \ref{mod:wdt} \textit{\nameref{mod:wdt}}\\
        \hline
        \hexprefix{}11 & MWDT1 CPU 复位 & CPU 复位 & 详见章节 \ref{mod:wdt} \textit{\nameref{mod:wdt}}\\
        \hline

    \end{longtable}
    注意：
    \begin{itemize}
        \item 芯片复位的触发源包括以下三项：
            \begin{itemize}
                \item 芯片上电
                \item 欠压检测器触发的芯片复位
                \item 超级看门狗 (SWD) 触发的芯片复位
            \end{itemize}
        \item 欠压检测器在检测到欠压状态时，将根据寄存器配置，选择触发系统复位或者芯片复位。详见章节 \ref{mod:lowpowm} \textit{\nameref{mod:lowpowm}}。
    \end{itemize}

\section{系统时钟}

\subsection{概述}

ESP32-S2 提供了多种不同频率的时钟选择，可以灵活配置 CPU、外设、以及 RTC 的工作频率，以满足不同功耗和性能需求。下图 \ref{figure:Clock} 为系统时钟结构。

\begin{figure}[H]
    \centering
    \includegraphics[width=0.9\textwidth]{\modulefiles/figures/sys_clk_arch.eps}
    \caption{系统时钟}
    \label{figure:Clock}
\end{figure}

\subsection{时钟源}
ESP32-S2 有三种时钟源：外部晶振、内部 PLL 和震荡电路，用于生成多种时钟。这些时钟根据频率不同可分为以下三种类型。
\begin{itemize}
    \item 快速时钟：供 CPU 及数字外设等高速设备使用
        \begin{itemize}
            \item PLL\_CLK：320 MHz 或 480 MHz 内部 PLL 时钟；
            \item XTAL\_CLK：40 MHz 外部晶振时钟。
        \end{itemize}
    \item 低功耗慢速时钟：供低功耗设备使用，包括功耗管理单元以及低功耗外设等
        \begin{itemize}
            \item XTAL32K\_CLK：32 kHz 外部晶振时钟；
	        \item RC\_FAST\_CLK：内置 8 MHz 振荡器的 N 分频时钟，频率可调；
	        \item RC\_FAST\_DIV\_CLK：由 RC\_FAST\_CLK 经 256 分频所得，频率为 RC\_FAST\_CLK / 256。当 RC\_FAST\_CLK 的初始频率为 8 MHz 时，该时钟以 31.250 kHz 的频率运行；
	        \item RC\_SLOW\_CLK：90 kHz 内部低功耗时钟，频率可调。
	    \end{itemize}
    \item 音频时钟：供音频相关设备使用
        \begin{itemize}
            \item APLL\_CLK：16 MHz \verb+~+ 128 MHz 内部 Audio PLL 时钟。
	    \end{itemize}
\end{itemize}

\subsection{CPU 时钟}
\label{sec:cpu-clock}
如图 \ref{figure:Clock} 所示，CPU\_CLK 为 CPU 主时钟，CPU 在最高效工作模式下，主频可以达到 240 MHz。同时，CPU 能够在超低频下工作（通常为2 MHz），以减少功耗。

CPU\_CLK 由 \hyperref[fielddesc:SYSTEMSOCCLKSEL]{SYSTEM\_SOC\_CLK\_SEL} 来选择时钟源，允许选择 PLL\_CLK、APLL\_CLK、RC\_FAST\_CLK 或 XTAL\_CLK 作为 CPU\_CLK 的时钟源。具体请参考表 \ref{table:cpu clksrc} 和表 \ref{table:cpu clk}。

\begin{tiplisting}
ESP32-S2 必须有外置主晶振时钟才可运行。
\end{tiplisting}

\begin{longtable}[c]{|p{8.5cm}|p{7.5cm}| }
       \caption{CPU\_CLK 源}
		\label{table:cpu clksrc} \\
	    \hline
        \rowcolor{lightgray}
        \hyperref[fielddesc:SYSTEMSOCCLKSEL]{SYSTEM\_SOC\_CLK\_SEL} 值 & 时钟源 \\
        \hline
        0 & XTAL\_CLK \\
        \hline
        1 & PLL\_CLK \\
        \hline
        2 & RC\_FAST\_CLK \\
        \hline
        3 & APLL\_CLK \\
        \hline
\end{longtable}

\begin{table}[hb]
	\caption{CPU\_CLK 选择}
    \label{table:cpu clk}
     \begin{tabular}[h]{ |p{3.3cm}|c|c|c|p{9.5cm}|}
        \hline
        \rowcolor{lightgray}
        时钟源 & SEL\_0* & SEL\_1* & SEL\_2* & CPU 时钟频率\\
        \hline
		\multirow{2}{*} {XTAL\_CLK} & \multirow{2}{*} {0} & \multirow{2}{*} {-} & \multirow{2}{*} {-} & CPU\_CLK = XTAL\_CLK / (\hyperref[fielddesc:SYSTEMPREDIVCNT]{SYSTEM\_PRE\_DIV\_CNT} + 1) \newline
         		\hyperref[fielddesc:SYSTEMPREDIVCNT]{SYSTEM\_PRE\_DIV\_CNT} 默认值为 1，范围 0 \verb+~+ 1023。\\
        \hline
       \multirow{2}{*} {PLL\_CLK (480 MHz)} & \multirow{2}{*} {1} & \multirow{2}{*} {1} & \multirow{2}{*} {0} & CPU\_CLK = PLL\_CLK / 6 \newline
         		CPU\_CLK 频率为 80 MHz。 \\
        \hline
       \multirow{2}{*} {PLL\_CLK (480 MHz)} & \multirow{2}{*} {1} &\multirow{2}{*} {1} & \multirow{2}{*} {1}  & CPU\_CLK = PLL\_CLK / 3 \newline
       			CPU\_CLK 频率为 160 MHz。 \\
        \hline
       \multirow{2}{*} {PLL\_CLK (480 MHz)} & \multirow{2}{*} {1} & \multirow{2}{*} {1} & \multirow{2}{*} {2}  & CPU\_CLK = PLL\_CLK / 2 \newline
       			CPU\_CLK 频率为 240 MHz。 \\
        \hline
        \multirow{2}{*} {PLL\_CLK (320 MHz)}  & \multirow{2}{*} {1}  & \multirow{2}{*} {0} & \multirow{2}{*} {0}  & CPU\_CLK = PLL\_CLK / 4 \newline
       			CPU\_CLK 频率为 80 MHz。 \\
        \hline
        \multirow{2}{*} {PLL\_CLK (320 MHz)}  & \multirow{2}{*} {1}  & \multirow{2}{*} {0} & \multirow{2}{*} {1}  & CPU\_CLK = PLL\_CLK / 2 \newline
       			CPU\_CLK 频率为 160 MHz。 \\
        \hline
       \multirow{2}{*} {RC\_FAST\_CLK}  & \multirow{2}{*} {2}  & \multirow{2}{*} {-}  & \multirow{2}{*} {-} & CPU\_CLK = RC\_FAST\_CLK / (\hyperref[fielddesc:SYSTEMPREDIVCNT]{SYSTEM\_PRE\_DIV\_CNT} + 1) \newline
        		\hyperref[fielddesc:SYSTEMPREDIVCNT]{SYSTEM\_PRE\_DIV\_CNT} 默认值为 1，范围 0 \verb+~+ 1023。\\
        \hline
        APLL\_CLK & 3 & 0 & 0 & CPU\_CLK = APLL\_CLK / 4 \\
        \hline
        APLL\_CLK & 3 & 0 & 1 & CPU\_CLK = APLL\_CLK / 2 \\
        \hline
    \end{tabular}
    *SEL\_0: 寄存器 \hyperref[fielddesc:SYSTEMSOCCLKSEL]{SYSTEM\_SOC\_CLK\_SEL} 的值； \\
    *SEL\_1: 寄存器 \hyperref[fielddesc:SYSTEMPLLFREQSEL]{SYSTEM\_PLL\_FREQ\_SEL} 的值； \\
    *SEL\_2: 寄存器 \hyperref[fielddesc:SYSTEMCPUPERIODSEL]{SYSTEM\_CPUPERIOD\_SEL} 的值。
    \end{table}

注意：
    \begin{itemize}
        \item 当 CPU 选择 XTAL\_CLK 用作时钟源，通过配置寄存器 \hyperref[fielddesc:SYSTEMPREDIVCNT]{SYSTEM\_PRE\_DIV\_CNT} 调节 XTAL\_CLK 的时钟分频系数时，需要遵循以下规则：
        \begin{itemize}
            \item 目标分频系数为 x 分频（x 不等于 1 分频 ）且当前分频系数为 2 分频时（\hyperref[fielddesc:SYSTEMPREDIVCNT]{SYSTEM\_PRE\_DIV\_CNT} = 1），需要先将分频系数配置为 1 分频（\hyperref[fielddesc:SYSTEMPREDIVCNT]{SYSTEM\_PRE\_DIV\_CNT} = 0），再调节分频系数为 x 分频（\hyperref[fielddesc:SYSTEMPREDIVCNT]{SYSTEM\_PRE\_DIV\_CNT} = x - 1）；
            \item 目标分频系数为 2 分频且当前分频系数为 x 分频（\hyperref[fielddesc:SYSTEMPREDIVCNT]{SYSTEM\_PRE\_DIV\_CNT} = x - 1），需要先将分频系数调节为 1 分频（\hyperref[fielddesc:SYSTEMPREDIVCNT]{SYSTEM\_PRE\_DIV\_CNT} = 0），再调节分频系数为 2 分频（\hyperref[fielddesc:SYSTEMPREDIVCNT]{SYSTEM\_PRE\_DIV\_CNT} = 1）；
            \item 如需调至其它目标分频系数 x，直接调整分频系数即可（\hyperref[fielddesc:SYSTEMPREDIVCNT]{SYSTEM\_PRE\_DIV\_CNT} = x - 1）。
        \end{itemize}
    \end{itemize}

\subsection{外设时钟}
外设所需要的时钟包括 APB\_CLK、REF\_TICK、LEDC\_PWM\_CLK、APLL\_CLK 和 PLL\_F160M\_CLK。\\
下表 \ref{table:peri clk owner} 为接入各个外设的时钟。

\begin{longtable}[c]{|p{2.5cm}|c|c|c|c|c| }
 \caption{外设时钟用法}
    \label{table:peri clk owner} \\
  \hline
        \rowcolor{lightgray}
        外设 & APB\_CLK & REF\_TICK & LEDC\_PWM\_CLK & APLL\_CLK & PLL\_F160M\_CLK\\
        \hline
        TIMG & Y & Y &  &  &  \\
        \hline
        I2S & Y &  &  & Y & Y \\
        \hline
        UHCI & Y &  &  &  &  \\
        \hline
        UART & Y & Y &  &  &  \\
        \hline
        RMT & Y & Y &  &  &  \\
        \hline
        LED\_PWM & Y & Y & Y &  &  \\
        \hline
        I2C & Y & Y &  &  &  \\
        \hline
        SPI & Y &  &  &  &  \\
        \hline
        PCNT & Y &  &  &  &  \\
        \hline
        eFuse Controller & Y &  &  &  &  \\
        \hline
        SARADC/DAC & Y &  &  & Y &  \\
        \hline
        USB & Y &  &  &  &  \\
        \hline
        CRYPTO &  &  &  &  & Y \\
        \hline
        TWAI Controller & Y &  &  &  &  \\
        \hline
        System Timer & Y &  &  &  &  \\
        \hline
\end{longtable}

\subsubsection{APB\_CLK 源}

如表 \ref{table:apb_clk_src} 所示，APB\_CLK 的频率由 CPU\_CLK 的时钟源决定。

 \begin{longtable}[c]{ |p{7.5cm}|p{7.5cm}| }
      \caption{APB\_CLK 源}
    \label{table:apb_clk_src} \\
        \hline
        \rowcolor{lightgray}
        CPU\_CLK 源 & APB\_CLK频率 \\
        \hline
        \endfirsthead
        \hline
        \rowcolor{lightgray}
        CPU\_CLK 源 & APB\_CLK \\
        \hline
        \endhead
        \endfoot
        PLL\_CLK & 80 MHz \\
        \hline
        APLL\_CLK & CPU\_CLK / 2 \\
        \hline
        XTAL\_CLK & CPU\_CLK \\
        \hline
        RC\_FAST\_CLK & CPU\_CLK \\
        \hline
\end{longtable}

\newpage
\subsubsection{REF\_TICK}\label{sec:ref_tick}

REF\_TICK 时钟频率由 APB\_CLK 分频产生，APB\_CLK 时钟频率由 CPU\_CLK 源决定。REF\_TICK 的时钟频率应固定，因此当切换 CPU\_CLK 源时，应配置分频寄存器，使其频率固定。分频寄存器如何配置见表 \ref{table:ref_tick}：

\begin{longtable}[c]{|l|l|l|}
      \caption{REF\_TICK}
    \label{table:ref_tick}
        \\\hline
        \rowcolor{lightgray}
        CPU\_CLK 源 & APB\_CLK 时钟频率 & REF\_TICK 频率\\
        \hline
        \endfirsthead
        \hline
        \rowcolor{lightgray}
        CPU\_CLK 源 & APB\_CLK 时钟频率 & REF\_TICK 频率\\
        \hline
        \endhead
        \endfoot
        PLL\_CLK & 80 MHz & APB\_CLK/(\hyperref[regdesc:SYSCONTICKCONFREG]{SYSCON\_XTAL\_TICK\_NUM} + 1)\\
        \hline
        APLL\_CLK & CPU\_CLK/2 & APB\_CLK/(\hyperref[regdesc:SYSCONTICKCONFREG]{SYSCON\_XTAL\_TICK\_NUM} + 1)\\
        \hline
         XTAL\_CLK & CPU\_CLK & APB\_CLK/ (\hyperref[regdesc:SYSCONTICKCONFREG]{SYSCON\_XTAL\_TICK\_NUM} + 1)\\
        \hline
        RC\_FAST\_CLK & CPU\_CLK & APB\_CLK/(\hyperref[regdesc:SYSCONTICKCONFREG]{SYSCON\_CK8M\_TICK\_NUM} + 1)\\
        \hline
\end{longtable}

例如，当 REF\_TICK 时钟频率固定为 1 MHz 时，如果 CPU\_CLK 源为 PLL\_CLK，则 REF\_TICK 频率 = 80 MHz/(\hyperref[regdesc:SYSCONTICKCONFREG]{SYSCON\_XTAL\_TICK\_NUM} + 1) = 1 MHz，那么 \hyperref[regdesc:SYSCONTICKCONFREG]{SYSCON\_XTAL\_TICK\_NUM} 应配置为 79 (\hexprefix{}4F)。



 \subsubsection{LEDC\_PWM\_CLK 源}

LEDC\_PWM\_CLK 时钟源由寄存器 \hyperref[fielddesc:LEDCAPBCLKSEL]{LEDC\_APB\_CLK\_SEL} 决定，如表 \ref{table:led_sclk} 所示。

    \begin{longtable}[c]{ |p{7.5cm}|p{7.5cm}| } \caption{LEDC\_PWM\_CLK 源}
    \label{table:led_sclk} \\
        \hline
        \rowcolor{lightgray}
        \hyperref[fielddesc:LEDCAPBCLKSEL]{LEDC\_APB\_CLK\_SEL} 值 & LEDC\_PWM\_CLK 源  \\
        \hline
        0（默认）& - \\
        \hline
        1 & APB\_CLK \\
        \hline
        2 & RC\_FAST\_CLK \\
        \hline
        3 & XTAL\_CLK \\
        \hline
\end{longtable}

\subsubsection{APLL\_SCLK 源}

APLL\_CLK 来自内部 PLL\_CLK，其输出频率通过使用 APLL 配置寄存器来配置。配置方式见第 \ref{audiopll} 节。

\subsubsection{PLL\_F160M\_CLK 源}

PLL\_F160M\_CLK 是 PLL\_CLK 根据当前 PLL 的频率分频所得。


\subsubsection{时钟源注意事项}

需要与其他时钟配合工作的外设（如 RMT、I2C 等）一般在选择 PLL\_CLK 时钟源的情况下工作。若频率发生变化，外设需要通过修改配置才能以同样的频率工作。接入 REF\_TICK 的外设允许在切换时钟源的情况下，不修改外设配置即可工作。详情请参考表 \ref{table:peri clk owner}。

LED 模块能将 RC\_FAST\_CLK 作为时钟源使用，即在 APB\_CLK 关闭的时候，LED 也可工作。换而言之，当系统处于低功耗模式时，
其他外设都将停止工作（APB\_CLK 关闭），但是 LED 仍然可以通过 RC\_FAST\_CLK 来正常工作。



\subsection{Wi-Fi 时钟}
Wi-Fi 必须在 APB\_CLK 时钟源选择 PLL\_CLK 下才能工作。只有当 Wi-Fi 进入低功耗模式时，才能暂时关闭 PLL\_CLK。

LOW\_POWER\_CLK 允许选择 XTAL32K\_CLK、XTAL\_CLK、RC\_FAST\_CLK、RTC\_SLOW\_CLK（RTC 当前所选的慢速时钟）用于 Wi-Fi 的低功耗模式。


\subsection{RTC 时钟}

\hypertarget{SLOW_CLK}{RTC\_SLOW\_CLK} 和 RTC\_FAST\_CLK 的时钟源为低频时钟。RTC 模块能够在大多数时钟源关闭的状态下工作。


RTC\_SLOW\_CLK 允许选择 RC\_SLOW\_CLK、XTAL32K\_CLK 或 RC\_FAST\_DIV\_CLK，用于驱动功耗模块。

RTC\_FAST\_CLK 允许选择 XTAL\_CLK 的分频时钟或 RC\_FAST\_CLK，用于驱动片上传感器模块。

\subsection{音频 PLL 时钟}\label{audiopll}
音频应用和其他对于数据传输时效性要求很高的应用都需要高度可配置、低抖动并且精确的时钟源。来自系统时钟的时钟源可能会携带抖动，并且不支持高精度的时钟频率配置。

在满足应用需求的前提下，为了最大限度地降低系统成本，ESP32-S2 集成了专门用于 I2S 外设的音频 PLL。用户可使用 APLL 时钟对 I2S 模块进行计时。

Audio PLL 公式如下：
$$f_{\rm{out}} = \frac{f_{\rm{xtal}}(\rm{sdm2}+\frac{\rm{sdm1}}{2^8}+\frac{\rm{sdm0}}{2^{16}}+4)}{2(odiv+2)} $$

其中，
\begin{itemize}
    \item \begin{math}f_{\rm{xtal}}\end{math}：晶振频率，通常为 40 MHz
    \item sdm0：可配参数 0 \verb+~+ 255
    \item sdm1：可配参数 0 \verb+~+ 255
    \item sdm2：可配参数 0 \verb+~+ 63
    \item odiv：可配参数 0 \verb+~+ 31
    \item 公式的分子频率工作范围在 350 MHz \verb+~+ 500 MHz
$$ 350MHz<{f_{\rm{xtal}}(\rm{sdm2}+\frac{\rm{sdm1}}{2^8}+\frac{\rm{sdm0}}{2^{16}}+4)}<500MHz $$
\end{itemize}

Audio PLL 可通过使能寄存器 \hyperref[fielddesc:RTCCNTLPLLAFORCEPU]{RTC\_CNTL\_PLLA\_FORCE\_PU} 强行打开，或者使能寄存器 \hyperref[fielddesc:RTCCNTLPLLAFORCEPD]{RTC\_CNTL\_PLLA\_FORCE\_PD} 强行关闭，关闭优先级大于打开优先级。当 \hyperref[fielddesc:RTCCNTLPLLAFORCEPU]{RTC\_CNTL\_PLLA\_FORCE\_PU} 和 \hyperref[fielddesc:RTCCNTLPLLAFORCEPD]{RTC\_CNTL\_PLLA\_FORCE\_PD} 同时为 0 时，PLL 会跟随系统状态，当系统进入睡眠模式时自动关闭，系统被唤醒时自动打开。

\section{寄存器列表}\label{sec:resclk-reg-summ}

本小节的所有地址均为相对于 SYSCON 基地址的地址偏移量（相对地址），具体基地址见章节 \ref{mod:sysmem} \textit{\nameref{mod:sysmem}} 中的表 \ref{table:Peripheral_Address_Mapping} \textit{外设地址映射}。

\subfile{\modulefiles/apb-ctrl-summary_cn.tex}

\section{寄存器}\label{sec:resclk-reg-desc}

本小节的所有地址均为相对于 SYSCON 基地址的地址偏移量（相对地址），具体基地址见章节 \ref{mod:sysmem} \textit{\nameref{mod:sysmem}} 中的表 \ref{table:Peripheral_Address_Mapping} \textit{外设地址映射}。

\subfile{\modulefiles/apb-ctrl-reg_cn.tex}


\end{document}
